\section{The End of Thinking}

\begin{quotex}
Reading, for a man devoid of prior understanding, is like a blind man's looking in a mirror.
\flright{\textsc{Garuda Purana}}
\end{quotex}


\begin{quotex}
There is no necessary connection of literacy with culture.
\flright{\textsc{Ananda Coomaraswamy}, \emph{The Bugbear of Literacy}}
\end{quotex}


\paragraph{Oral Culture}


As we pointed out in The Great Migration\footnote{\url{https://gornahoor.net/?p=1393}}, tradition used to be passed on through the arts, particularly poetry and visual arts. During festivals, revealed texts, myths, and epic poems would be recited from memory to an audience able to concentrate for hours, far exceeding the abilities of the ADD minds of our contemporaries. Such festivals would unite the castes in a common worldview. With the rise of the printing press and literacy, oral culture has been lost. The printing press allegedly overcame the miasma and superstitions of the Middle Ages. Actually, the net effect was to create more viral and ineradicable myths. As \textbf{Jacques Ellul} points out, the purpose of mass education is to facilitate the flow of propaganda.

\paragraph{Scholasticism}


At the peak of Scholasticism, when a Traditional metaphysics based on spiritual experience was being developed, literacy was restricted to the few. Expertise in logic and reasoning was a prerequisite. In such an atmosphere of high standards and minimal interference, thinkers could proceed in relative peace. As Steiner describes it:

\begin{quotex}
We have only to remember, of course, in what surroundings this thinking took place. It did not take place in the midst of the noisy world as it does today. It took place in the cloister or somewhere far from the world. It was a thinking that was absorbed in the life of thought and was also able, through other circumstances, to arrive at a pure thought-technique. As a matter of fact it is very difficult today to build up such a pure activity of thought. For no sooner do we seek to make public any such thought-activity, in which the sole aim is to allow thought to follow thought in accordance with their content, than \emph{along come stupid people, raising every conceivable objection and interposing their violent prejudices}. In such circumstances that inner quiet is very soon lost to which the thinkers of the twelfth and thirteenth centuries were able to devote themselves. \textbf{For in the life of their day they were not compelled to pay such disproportionate attention to the opposition of the uneducated}.
\flright{\textsc{Rudolf Steiner}, \emph{The Redemption of Thinking}}
\end{quotex}


\paragraph{The Matrix}


The modern mind can only see in Scholasticism a dull abstraction, rather than a living and direct experience with the world of ideas. As the experience of, and hence the reality of, ideas began to fade, the position called Nominalism began to gain adherents. Instead of grasping the essence of things, the nominalist understands ideas as mere words to categorize similar experiences.

Galileo denied secondary qualities and situated qualities such as colour in the mind alone. Hume even denied the existence of causality. Kant took this line of thought to its complete logical conclusion. We cannot know ``reality'' as the thing-in-itself, but only our subjective impositions and interpretations. In effect, we live in the ``Matrix''. We undergo a series of subjective experiences, yet are always unsure about the relationship between those experiences and some other ``external'' world that may be creating them.

\paragraph{Hostile Environment}


Kantian subjectivism is the dominant way of thinking of moderns, and even that of many who oppose the modern world. If subjectivity is the only ``reality'' that we can know, then contemporary trends become easy to understand. Legal ideas such as a hostile environment, wherein the mere subjective impression of intimidation, harassment, oppression, ``the gaze'', and so on, make perfect sense within that framework. Those who propose, in its stead, objective standards or the ``defense of truth'', are seen as irrelevant at best, or else, at worst, as proposing another tool of oppression. Furthermore, these latter are often Kantians themselves, so their objections are hollow.

\paragraph{The Helplessness of Science}


There are those who appeal to science, in the name of ``realism'', to oppose modern trends. However, traditional sciences used to be the study of ``essences'' or the ``whatness'' of things, while modern science restricts itself to the interrelationships among phenomena. A scientific hypothesis becomes a creative work of literature, of a specified type and requirements. While a sonnet consists of 14 lines with a specific rhyming pattern, a hypothesis must ``save the appearances'', that is, account for observed phenomena. Unfortunately, when it comes to human studies, there is no way to keep variables constant and perform experiences. Hence, other competing, and mutually contradictory, hypotheses can be proposed with no absolute way to arbitrate between them. Science is helpless in such situations.

\paragraph{Conclusion}


Those rebels against the modern world need to return to a more qualitative understanding of intelligence, apart from IQ. The Kantian mindset much be overcome which will require a conscious effort. Intelligence is the grasping of essences, the world of ideas, which is more like a ``seeing''. Proper thinking is the awareness of spiritual reality, not a subjective process.


\flrightit{Posted on 2011-05-18 by Cologero}

\begin{center}* * *\end{center}

\begin{footnotesize}
\begin{sffamily}
\texttt{White Rabbit on 2015-12-25 at 04:35 said:}

IQ doesn't even necessarily translate to rationality. The cleverer someone is, the easier they can defend stupidity and lies. Is that kind of mind ``intelligent''?

\end{sffamily}
\end{footnotesize}

